
\subsecao{Contextualization}

Social science has long suffered from a singular, fatal constraint: the inability to run controlled experiments on history itself. We cannot collapse an economy just to test a hypothesis on inflation, nor can we deploy a radical social policy merely to observe the fallout. Synthetic societies dismantle this barrier. By orchestrating thousands of autonomous, cognitively complex agents into a living \textit{in silico} mirror of our world, this technology grants the power to simulate the consequences of an intervention before it happens — transforming the speculative art of forecasting into an experimental science, a "sandboxed reality" where stakeholders, from central banks to policy designers, can stress-test their most volatile interventions against a population that negotiates, adapts, and evolves, without risking a single real-world consequence.

Within this scenario, Large Language Models (LLMs) transformed the field by overcoming the historical limitation of rigid agent reasoning. LLMs opened the possibility for agents to exhibit flexible, human-aligned reasoning patterns, directly enabling synthetic societies capable of complex emergent behavior. This shift is central to the present research agenda, because it explains why LLM-based agents can serve as cognitively rich components in large-scale social simulations.

Even before the emergence of LLM-based agents, attempts to simulate social dynamics relied primarily on traditional Agent-Based Modeling (ABM). Classical artificial societies captured important macropatterns — such as segregation, mobility, and economic exchange — through rule-based agents. However, these agents operated under rigid, predefined behavioral rules and simplified utility functions, which prevented them from representing the cognitive richness, heterogeneity, and adaptive reasoning observed in real populations. This historical limitation establishes the conceptual inflection point that motivates the transition toward cognitively enhanced, LLM-driven agents.

These advances allow the construction of synthetic societies populated by LLM agents displaying heterogeneous backgrounds, beliefs, and economic profiles. Researchers can now build digital laboratories to test economic policies, model social conflicts, and study collective decision-making with greater realism, capturing deeper heterogeneity in systemic outcomes.

However, a major scientific gap remains regarding calibration and structural validity. LLM-agent simulations are often developed with a focus on narrative plausibility rather than statistical fidelity. As Horton (2023) highlights, while LLMs can act as simulated economic agents, their utility for the social sciences depends entirely on their ability to replicate structural empirical regularities rather than mere linguistic coherence. The field currently relies on ad hoc prompt designs that lack standardized mechanisms for directly injecting real-world census data or socioeconomic distributions into the agent-generation pipeline. Without a rigorous methodology to ground these agents in empirical data — transforming them from "stochastic parrots" into calibrated personas — simulation outcomes risk becoming hallucinatory artifacts rather than reproducible social insights. This research aims to bridge this gap.

This monograph documents a specific instance of that calibration failure that is invisible in single-agent settings but produces pathological outcomes in multi-agent environments: \textbf{RLHF-aligned LLMs are individually virtuous but collectively pathological.} Across multiple social-dilemma game types — public goods games, prisoner's dilemmas, stag hunts, and ultimatum games — LLMs fine-tuned via Reinforcement Learning from Human Feedback (RLHF; Ouyang et al., 2022) exhibit cooperation rates far exceeding both the Nash equilibrium and empirically observed human behavior. We call this the \textbf{RLHF Cooperative Bias} and formalize it via the index \texttt{B\_RLHF = TV(π, π\_uniform)}.

The mechanism is structural. RLHF trains models on single-agent human-preference data in which the evaluator is always cooperative and well-intentioned — a context in which helpfulness and cooperation are synonymous. This produces a strong cooperative prior that overgeneralizes to multi-agent environments: placed in a social dilemma alongside agents with competing interests, the RLHF-tuned model cooperates with every agent as if each were the RLHF evaluator. It has no learned representation of adversarial interaction, no basis for trust discrimination, and no training signal indicating that cooperation can be individually costly. The result is a synthetic society resembling a utopia: frictionless cooperation, unnatural egalitarianism, and zero social friction — a world that has never existed. This is an alignment failure, not a simulation artifact, and as LLMs are increasingly deployed in multi-agent contexts (AI councils, automated negotiation, agentic workflows with competing objectives) measuring and mitigating this bias becomes a prerequisite for deployment.

\subsecao{Problem Delimitation}

This research addresses the absence of a robust, empirically grounded framework capable of using demographic datasets, such as census or survey data, together with a reproducible pipeline for building, calibrating, and testing LLM-based synthetic societies. Although the technology to construct LLM agents already exists, there is still no established method to ensure that their collective behavior aligns with real-world socioeconomic patterns, to systematically translate demographic datasets into agent personas, or to validate whether the aggregate behavior of these agents reflects observed empirical regularities.

Moreover, existing studies rarely distinguish narrative plausibility from structural fidelity, making it difficult to determine whether simulated outcomes emerge from meaningful social mechanisms or from uncontrolled linguistic artifacts inherent to LLMs. The lack of standardized grounding procedures also undermines reproducibility, as current approaches depend on ad hoc prompt engineering and non-transparent calibration heuristics.

Within this context, this work delimits its scope to a fundamental methodological problem: creating agent personas directly from real data, defining consistent interaction rules, and establishing metrics that measure behavioral accuracy. The goal is to articulate the scientific standards required for valid applications and to operationalize them in a concrete, reproducible system — the Behavioral Grounding Framework — whose collective outputs can be quantitatively compared against empirical socioeconomic baselines.

\subsecao{Justification}

The lack of a rigorous method threatens the scientific legitimacy of LLM-based computational social science. If policy recommendations derived from synthetic societies rely on uncalibrated agents that do not reflect actual income distributions, policymakers may implement interventions that are ineffective or even detrimental when applied to real populations. Without standardized grounding in real data — income distributions, educational levels, regional demographics — findings from synthetic societies cannot be replicated and may lead to misleading interpretations, limiting their value for research and policy.

The relevance of this work is grounded in advancing synthetic societies from exploratory demonstrations to empirically calibrated simulations capable of supporting analytical studies. Treating agents not as fictional characters but as statistical representations of real demographic groups provides the necessary foundation for structured analyses of collective behavior and socioeconomic dynamics. Researchers increasingly require methods that move beyond ad hoc prompt engineering toward reproducible and transparent procedures; likewise, public institutions that consider using synthetic societies for forecasting, evaluation of economic measures, or scenario testing need models supported by measurable fidelity rather than narrative coherence alone.

This research is timely because current approaches still lack methodological consistency and clear validation mechanisms. By focusing on how to ground agent construction and assess collective behavior against real-world baselines, this study establishes a conceptual and operational foundation that enables empirical implementation while maintaining rigor and reproducibility.

\subsecao{Objectives}

This work has \textbf{two explicit, pre-registered aims}, both pursued throughout: (Aim 1 — measurement) to formalize quantitative metrics for detecting RLHF-induced behavioral distortion and for assessing behavioral fidelity in LLM-based synthetic societies; and (Aim 2 — propagation test) to use those metrics to test whether empirical grounding in real socioeconomic microdata propagates through LLM agents into realistic collective behavior. The metrics work (Aim 1) is therefore a planned, first-class contribution, not a by-product; the grounding propagation test (Aim 2) is the planned empirical hypothesis whose outcome — reported honestly whether positive or null — constitutes the central finding.

\textbf{General objective.} To formalize and apply quantitative metrics of RLHF behavioral distortion and behavioral fidelity within an empirically grounded LLM-based synthetic-society framework, and to test whether socioeconomic grounding propagates through the LLM decision layer into realistic collective behavior.

\textbf{Specific objectives.}

\begin{itemize}
\item (Aim 1) To formalize statistical distance measures and behavioral-consistency criteria — the RLHF Bias Index \texttt{B\_RLHF = TV(π, π\_uniform)} and the Behavioral Realism Metric (BRM) — that detect alignment-induced distortion and quantify the fidelity of synthetic societies against real social and economic patterns;
\item To systematize the conceptual foundations required to transform demographic and socioeconomic datasets into semantically coherent LLM-based agent personas, and to define the architectural principles and behavioral constraints structuring the proposed Behavioral Grounding Framework;
\item To synthesize a reproducible pipeline integrating data grounding, persona construction, interaction modeling, and fidelity assessment, validated empirically against real-world socioeconomic baselines (ESS, Eurostat, WVS, and an independent public-goods-game benchmark);
\item (Aim 2) To test, under a pre-registered hypothesis battery, whether empirical grounding propagates through the LLM decision policy into measurable behavioral differentiation, and to characterize the outcome with the metrics of Aim 1.
\end{itemize}

These objectives are operationalized as seven research questions:

\begin{itemize}
\item \textbf{RQ1 (Primary alignment finding):} Is the RLHF cooperative bias universal across social-dilemma game types, or specific to the public-goods setting?
\item \textbf{RQ2 (Grounding efficacy):} Does ESS grounding significantly mitigate \texttt{B\_RLHF} by providing trust- and risk-calibrated priors that override the RLHF cooperative default?
\item \textbf{RQ3 (Realism):} Can ESS-grounded LLM agents reproduce the macroeconomic and network-topological phenomena of real human populations?
\item \textbf{RQ4 (Cross-model scope):} Is the bias universal across RLHF alignment families, or moderated by alignment methodology?
\item \textbf{RQ5 (Memory):} What is the independent contribution of each memory tier to persona fidelity and behavioral consistency?
\item \textbf{RQ6 (Cross-cultural):} Does the grounding function recover cross-cultural behavioral variation measured by ESS and independently validated by WVS?
\item \textbf{RQ7 (Robustness):} Are grounded societies resilient to adversarial perturbations, economic shocks, and network-topology variation?
\end{itemize}

\subsecao{Methodological Approach}

This work adopts a constructive, theoretically driven, and empirically validated methodological approach, structured in four integrated movements. First, a focused literature review is conducted on agent-based modeling, cognitive and generative architectures, synthetic populations, and validation metrics, in order to delineate the conceptual and methodological limitations of existing approaches and extract the design requirements for the framework. Second, the \textbf{Behavioral Grounding Framework} is proposed, formally specifying the stages of data grounding, persona construction, interaction modeling, and the definition of quantitative indicators of fidelity and inequality, always anchored in empirical socioeconomic distributions. Third, the framework is \textbf{implemented} as a reproducible software artifact with deterministic seeding, dual SQL/Graph retrieval, hierarchical memory, and a cryptographic reproducibility witness. Fourth, synthetic societies generated under the framework are \textbf{empirically evaluated} through distributional distance measures and inequality indices against real-world baselines (ESS, Eurostat, WVS, and an independent public-goods-game benchmark), under a pre-registered hypothesis battery with family-wise error correction.

Whereas the project phase of this Trabalho de Conclusão de Curso defined the conceptual route (theoretical review → conceptual formulation → validation strategy), the present monograph carries that route through to execution, reporting the empirical results obtained from the implemented framework.

\subsecao{Structure of the Work}

This monograph is organized into seven chapters. After this Introduction, Chapter 2 reviews the literature, develops the theoretical framework, surveys related works, and presents the validation framework and gap analysis. Chapter 3 formalizes the Behavioral Grounding Framework methodology: the formal tuple, the fidelity metrics and their formal results, the system architecture, decision policies, prompt engineering and ablation design, anti-drift engineering, the causal-identification strategy, the experimental conditions and hypothesis pre-registration, and the software artifact. Chapter 4 details the experimental setup (data, model, ethics, statistical power, stress tests, phase-transition sweeps, and cross-model validation). Chapter 5 reports the results, from the proof of concept through the macroeconomic, topological, robustness, cross-cultural, cross-model, memory-ablation, and mechanism analyses. Chapter 6 discusses the findings — centrally the Φ/P\_LLM dissociation — and their implications. Chapter 7 presents the limitations, broader impacts, conclusions, contributions, and future work. The References close the document.

